\section{Literature Review}

The application of Artificial Intelligence (AI) and Machine Learning (ML) in educational systems has rapidly expanded over the past decade, shifting from basic statistical analysis to highly complex predictive modeling. This chapter extensively reviews the existing literature surrounding educational data mining, the mathematical evolution of multi-label \parencite{tsoumakas2007multi} classification algorithms, and the critical, emerging need for Explainable AI (XAI) in high-stakes guidance systems. By critically examining the current state of academic research, this chapter establishes the robust theoretical foundation required for transitioning from simple, isolated academic performance prediction to an intelligent, mathematically explainable, and multi-pathway career recommendation framework.

\subsection{Theoretical Framework: From Academic Performance to Career Pathway}

Traditionally, educational guidance has been viewed by both educators and sociologists as a highly linear, deterministic progression: a student's past academic performance dictates their immediate educational choices, which subsequently permanently restrict or open up future career paths. Educational Data Mining \parencite{romero2010educational} (EDM) emerged as a formalized discipline precisely to quantify and analyze these transitions. According to foundational work by \textcite{romero2010educational}, EDM focuses on developing mathematical, statistical, and computational methods to discover unique, actionable patterns in the massive amounts of data coming from educational environments.

The theoretical framework underpinning this thesis rejects the linear, single-destination model and instead proposes a continuous, intelligent, and multi-branching pathway:
\[
\text{Academic Performance (Input)} \xrightarrow{\text{ML Prediction}} \text{Educational Pathway} \xrightarrow{\text{Recommendation}} \text{Career Pathway}
\]

In this mathematical context, the input consists of numerical scores from standardized tests. Let $X$ represent the high-dimensional feature space. For the scope of this thesis, $X = [x_{math}, x_{phys}, x_{bio}, x_{chem}]^T$, representing the core STEM subjects required for health sciences. The traditional, outdated algorithmic approach is to map $X$ to a single target variable $Y$ (a specific, isolated career). 

However, human potential, cognitive aptitude, and professional capabilities are inherently not mutually exclusive. A student excelling simultaneously in Biology and Chemistry could successfully pursue Medicine, Pharmacy, or Nursing with equal theoretical probability. Therefore, the critical theoretical shift required in modern EDM is moving away from the restrictive mapping $f: X \rightarrow Y$ and toward a multi-dimensional mapping $f: X \rightarrow \{0, 1\}^K$, where $K$ represents the total number of possible, overlapping career paths \autocite{baker2014educational}. This framework forms the absolute core of the thesis.

\subsection{Chronological Evolution of Educational Guidance Systems}

The methodology for guiding students toward appropriate career pathways has undergone a radical transformation over the past century, evolving from entirely subjective human assessments to highly mathematical, data-driven automated systems. 

In the late 20th century, guidance was fundamentally reliant on psychometric testing and manual counselor evaluation. Tools such as the Holland Occupational Themes (RIASEC) provided a structured, albeit entirely manual, framework for categorizing student personalities and matching them to predefined career clusters. While foundational, these systems were static. They relied on a student's self-reported psychological state at a single point in time, completely ignoring the longitudinal, empirical reality of their hard academic performance in subjects like Mathematics and Physics.

By the early 2000s, the advent of basic database management systems allowed educational institutions to deploy "Expert Systems." These early computational models utilized hard-coded, rules-based logic (e.g., IF Math $> 80$ AND Biology $> 75$ THEN Recommend Pre-Med). However, expert systems suffered from a fatal flaw known as the "knowledge acquisition bottleneck." Because human experts had to manually define every single rule and threshold, the systems were incredibly brittle, entirely unable to scale, and failed to capture the non-linear, complex covariance between different academic disciplines that define true cognitive aptitude.

It was not until the mass digitization of national educational records (Educational Data Mining) in the 2010s that true Machine Learning could be applied to this domain, shifting the paradigm from manually written rules to autonomously discovered mathematical patterns.

\subsection{Machine Learning in Educational Data Mining}

Predictive modeling in education has a rich history of utilizing historical data to forecast student success, retention, or dropout rates. Early academic studies primarily relied on foundational statistical methods such as logistic regression and basic, singular decision trees. While these models are highly interpretable, meaning human operators can easily extract the exact mathematical coefficients or rules dictating how the model makes decisions, they fundamentally fail to capture the complex, highly non-linear interactions between different academic subjects that truly define student capability.

For example, a standard Decision Tree partitions the feature space recursively using orthogonal hyperplanes. If $x_{math}$ represents the Mathematics score and $x_{bio}$ represents the Biology score, the tree uses a simple splitting criterion (such as Gini Impurity) to create a rule such as:
\[
\text{If } x_{math} > 70 \text{ and } x_{bio} > 75 \implies \text{Recommend Medicine}
\]
While this is mathematically simple and logically easy to understand for a human counselor, this rigid, box-like structure often leads to "overfitting" in complex educational datasets. The model learns the exact statistical noise of the training data rather than the underlying true pedagogical pattern. To overcome this critical limitation, EDM researchers began heavily utilizing ensemble methods. Ensemble learning leverages the mathematical principle of the "wisdom of crowds," combining multiple weak, high-variance models to create a single strong, highly generalized predictive model \autocite{romero2010educational}.

\subsection{Single-Label vs. Multi-Label \parencite{tsoumakas2007multi} Classification}

To fully appreciate the methodology of this thesis, one must understand the stark mathematical differences between traditional single-label classification and multi-label \parencite{tsoumakas2007multi} classification.

In standard single-label classification, the objective is to learn a hypothesis function that assigns a given instance $x_i$ to exactly one class out of a set of mutually exclusive classes. Mathematically, if $L = \{l_1, l_2, \dots, l_K\}$ is a finite, predefined set of labels, a single-label model predicts a single output $y_i \in L$. The probabilities across all classes must sum exactly to 1:
\[
\sum_{k=1}^{K} P(y=l_k | x) = 1
\]

However, as established in the theoretical framework, health sector career guidance is inherently a \textbf{Multi-Label Classification} problem. A student possesses a complex skill profile that makes them suitable for multiple careers simultaneously. In multi-label \parencite{tsoumakas2007multi} classification, the model is tasked with predicting a subset of labels $Y \subseteq L$ for each instance. This means the output is no longer a single scalar, but rather a vector $Y = [y_1, y_2, \dots, y_K]$, where $y_k \in \{0, 1\}$ indicates the relevance or recommendation of the $k$-th career \autocite{tsoumakas2007multi}. In this paradigm, the probabilities for each individual label are calculated independently, and their sum is not constrained to 1.

\textcite{zhang2013review} comprehensively categorize multi-label \parencite{tsoumakas2007multi} learning algorithms into two primary, distinct approaches:
\begin{enumerate}
    \item \textbf{Problem Transformation Methods:} These methods mathematically convert the complex multi-label \parencite{tsoumakas2007multi} problem into one or more traditional single-label problems, allowing standard algorithms (like basic SVMs or Decision Trees) to be used without internal modification.
    \item \textbf{Algorithm Adaptation Methods:} These methods rewrite the core internal mathematics of existing algorithms (such as altering the entropy calculation of a Decision Tree) to allow them to directly handle multi-label \parencite{tsoumakas2007multi} data.
\end{enumerate}
This thesis focuses heavily on Problem Transformation Methods, specifically seeking to overcome the limitations of the most basic transformation: Binary Relevance.

\subsubsection{Binary Relevance and the Independence Assumption}

The simplest, most naive problem transformation method is \textbf{Binary Relevance (BR)}. The BR algorithm operates by training $K$ completely independent binary classifiers, one for each distinct label in the dataset. For a new instance $x$, the final multi-label \parencite{tsoumakas2007multi} prediction is simply the logical union of the positive labels predicted by all $K$ individual classifiers.

While computationally fast and mathematically simple to implement, BR suffers from a fatal theoretical flaw: it assumes that all labels are completely statistically independent. Mathematically, it assumes:
\[
P(y_1, y_2, \dots, y_K | x) = \prod_{k=1}^{K} P(y_k | x)
\]
In the specific context of guiding students into highly specialized health sector careers, this independence assumption spectacularly fails. The probability of a student being suited for "Nursing" is highly positively correlated with the probability of them being suited for "Midwifery." If an intelligent model predicts a high likelihood for one, it should mathematically leverage that newly discovered information to increase the predicted probability of the other. BR ignores this entirely.

\subsubsection{Classifier Chains for Label Correlation}

To theoretically and mathematically overcome the naive independence assumption of Binary Relevance, \textcite{read2011classifier} introduced the highly effective \textbf{Classifier Chains (CC)} method. CC elegantly transforms the multi-label \parencite{tsoumakas2007multi} problem into a sequential chain of binary classifiers. The critical innovation is that subsequent classifiers in the chain utilize the actual predictions of preceding classifiers as additional input features, thereby forcing the model to learn the correlations between labels.

Let the original input feature vector be $x$. For $K$ target labels, a sequence of $K$ classifiers $C_1, C_2, \dots, C_K$ is trained. The training process proceeds as follows:
\begin{itemize}
    \item Classifier $C_1$ is trained on the original features $x$ to predict label $y_1$.
    \item Classifier $C_2$ is trained on an extended feature space consisting of $x$ appended with the true value of $y_1$ (i.e., $[x, y_1]$) to predict $y_2$.
    \item Generalizing this, Classifier $C_k$ is trained on the feature space $[x, y_1, y_2, \dots, y_{k-1}]$ to predict $y_k$.
\end{itemize}

During the inference (prediction) phase on a new, unseen student, the true labels $y$ are completely unknown. Therefore, the predictions $\hat{y}$ from previous classifiers in the chain are passed down instead:
\[
\hat{y}_k = C_k(x, \hat{y}_1, \dots, \hat{y}_{k-1})
\]
This chaining mechanism fundamentally changes the probability calculation. Instead of independent probabilities, CC models the joint probability of the labelset using the chain rule of probability:
\[
P(y_1, y_2, \dots, y_K | x) = P(y_1 | x) \cdot P(y_2 | x, y_1) \cdot \dots \cdot P(y_K | x, y_1, \dots, y_{k-1})
\]
This allows the model to capture deep, complex correlations between different career paths, making it vastly more suitable for intelligent educational guidance than standard Binary Relevance.

\subsection{Review of Selected Base Algorithms}

Within the Classifier Chain \parencite{read2011classifier} framework, a "base algorithm" must be chosen to act as the individual binary classifier ($C_k$) at each link in the chain. The mathematical complexity and learning capability of the CC model depend entirely on this base algorithm. This thesis explores four powerful, mathematically distinct algorithms:

\subsubsection{Decision Trees (DT)}
A Decision Tree recursively splits the input data space based on the feature that maximizes a specific purity metric, typically Information Gain. If a dataset $D$ has $C$ classes, the Shannon entropy $H(D)$ is defined mathematically as:
\[
H(D) = - \sum_{i=1}^{C} p_i \log_2(p_i)
\]
where $p_i$ is the empirical proportion of samples belonging to class $i$. The Information Gain for a split on feature $A$ is the difference between the parent node's entropy and the weighted sum of the child nodes' entropies. While highly interpretable, singular decision trees are incredibly prone to high variance and severe overfitting when dealing with deep, continuous numerical data like exam scores.

\subsubsection{Random Forest (RF)}
To solve the high variance problem of single decision trees, \textcite{breiman2001random} introduced Random Forests. A Random Forest \parencite{breiman2001random} is a sophisticated ensemble algorithm that combines the concepts of "bagging" (bootstrap aggregating) and random feature selection. It builds hundreds of deep, independent decision trees in parallel. Each tree is trained on a random bootstrap sample of the data, and at each split in the tree, only a random subset of features is considered. The final prediction for a regression or probability task is made by averaging the outputs of all trees:
\[
\hat{f}_{RF}(x) = \frac{1}{B} \sum_{b=1}^{B} f_b(x)
\]
where $B$ is the total number of trees. This mathematical averaging heavily reduces the variance of the model without significantly increasing the bias, making RF exceptionally robust against overfitting.

\subsubsection{XGBoost (Extreme Gradient Boosting)}
XGBoost, developed by \textcite{chen2016xgboost}, represents the absolute state-of-the-art for structured, tabular data analysis. Unlike Random Forest \parencite{breiman2001random}, which builds deep trees independently in parallel, XGBoost \parencite{chen2016xgboost} builds shallow trees sequentially. This is the core concept of "boosting." Each new tree is specifically trained to focus on correcting the mathematical errors (the residuals or gradients) made by the combination of all previous trees. 

The objective function optimized at iteration $t$ contains a specific loss function $l$ (which measures how far the prediction is from the truth) and a regularization term $\Omega$ (which penalizes overly complex trees to prevent overfitting):
\[
\mathcal{L}^{(t)} = \sum_{i=1}^{n} l\left(y_i, \hat{y}_i^{(t-1)} + f_t(x_i)\right) + \Omega(f_t)
\]
where the regularization term is defined as:
\[
\Omega(f) = \gamma T + \frac{1}{2} \lambda \sum_{j=1}^{T} w_j^2
\]
Here, $T$ is the number of leaves in the tree, $w_j$ are the leaf weights, and $\gamma$ and $\lambda$ are hyperparameter constants. By utilizing a second-order Taylor expansion to approximate the loss function, XGBoost \parencite{chen2016xgboost} optimizes this objective with extreme computational efficiency, making it a highly dominant algorithm in modern data science competitions and enterprise applications.

\subsubsection{Multi-Layer Perceptron (MLP)}
An MLP is the foundational architecture for deep artificial neural networks. It consists of an input layer, one or more hidden layers, and an output layer. Neurons in one layer are densely connected to neurons in the next, with each connection governed by a mathematical weight $w_{ij}$. The output $z_j$ of a neuron $j$ is calculated by applying a non-linear activation function $\sigma$ to the weighted sum of all its inputs plus a bias term $b_j$:
\[
z_j = \sigma \left( \sum_{i} w_{ij} x_i + b_j \right)
\]
For modern networks, the activation function $\sigma$ is typically the Rectified Linear Unit (ReLU), defined simply as $\max(0, x)$. The network learns by utilizing the backpropagation algorithm to calculate the gradient of the loss function with respect to every weight, updating the weights using optimizers like Stochastic Gradient Descent (SGD) or Adam. While MLPs possess the theoretical capacity (via the Universal Approximation Theorem) to model any continuous non-linear relationship, their deeply layered nature makes them the ultimate "black box," rendering it nearly impossible for a human educator to understand exactly how a recommendation was reached without external XAI tools.

\subsubsection{Limitations of Traditional Algorithms (SVM and Naive Bayes)}
While the aforementioned tree-based ensembles (XGBoost, RF) and MLPs represent the cutting edge for this specific problem, much of the early literature in educational data mining heavily relied on Support Vector Machines (SVM) and Naive Bayes classifiers. It is mathematically critical to understand why these algorithms are inherently sub-optimal for multi-label \parencite{tsoumakas2007multi} health sector guidance based on raw exam scores.

\textbf{Naive Bayes (NB)} operates on the strict mathematical assumption of conditional independence between features given the class label:
\[
P(x_{math}, x_{bio}, x_{chem} | Y=1) = P(x_{math}|Y=1) \cdot P(x_{bio}|Y=1) \cdot P(x_{chem}|Y=1)
\]
In reality, a student's performance in Physics is highly correlated with their performance in Mathematics. The Naive Bayes algorithm completely ignores this deep covariance matrix, leading to severely biased probability estimates and poor ranking performance in multi-label \parencite{tsoumakas2007multi} scenarios.

\textbf{Support Vector Machines (SVM)}, conversely, attempt to find a hyper-plane in an N-dimensional space that distinctly classifies the data points. While mathematically elegant (often utilizing the "kernel trick" to map data into higher dimensions for linear separability), SVMs are computationally expensive and scale very poorly ($O(n^3)$) with large educational datasets. Furthermore, SVMs natively output distance metrics to a hyperplane, not true probabilities. To generate the necessary ranked probabilities for career recommendations, SVM outputs must be artificially squeezed through Platt scaling (a logistic regression model trained on the SVM scores), which introduces an additional layer of unnecessary approximation error compared to algorithms like Random Forest \parencite{breiman2001random} that natively calculate empirical probabilities based on leaf-node distributions.

\subsection{Explainable AI (XAI) and Counterfactuals}

The rapid deployment of advanced ML systems (like XGBoost \parencite{chen2016xgboost} and MLPs) in high-stakes socio-technical areas such as healthcare, criminal justice, and education demands absolute algorithmic transparency. A career guidance model that simply outputs a probability matrix indicating "Recommend Medicine: 92\%" without any justification is highly unlikely to be trusted or adopted by students, parents, or government ministries.

Explainable AI (XAI) seeks to penetrate the black box \parencite{wachter2017counterfactual}, providing mathematical visibility into the decision-making process of complex ensemble and deep learning models. Foundational work by \textcite{lundberg2017unified} proposed SHAP (SHapley Additive exPlanations), a rigorous game-theoretic approach that assigns a specific, mathematically guaranteed contribution value to each feature (e.g., Mathematics score vs. Biology score) for every specific prediction made by the model. 

However, knowing \textit{why} a decision was made is often insufficient for human users seeking guidance. Furthermore, \textcite{wachter2017counterfactual} introduced the powerful concept of \textbf{Counterfactual Explanations}. Instead of just explaining why a model rejected a student for a specific career, counterfactuals mathematically explain exactly what would need to change in the student's profile to receive the desired decision. 

As further discussed by \textcite{barocas2020hidden}, counterfactuals are highly actionable. Mathematically, finding a counterfactual \parencite{wachter2017counterfactual} $x^*$ for an original, factual input $x$ involves setting up an optimization problem to minimize the distance $d$ between them, subject to the strict constraint that the model's prediction output changes to the desired class:
\[
x^* = \arg\min_{x'} d(x,x') \quad \text{subject to } f(x') \neq f(x)
\]
In a practical educational context, this advanced mathematical formulation allows the guidance system to generate statements such as: \textit{"You were recommended for Nursing based on your current profile. However, if your Chemistry raw mark had been exactly 12 points higher, the algorithm would have shifted your primary recommendation to Pharmacy."} This level of actionable, personalized insight is entirely absent from current educational data mining applications.

\subsection{Ethical AI and Algorithmic Fairness in Education}

As AI systems move from academic research to physical deployment in national education infrastructure, the mathematical concept of "fairness" becomes a critical limitation that must be addressed. A model trained on historical data is inherently at risk of learning and perpetuating historical biases \autocite{barocas2016big}. 

If, historically, female students were disproportionately discouraged from pursuing advanced degrees in STEM, a naive machine learning algorithm will observe that correlation in the training data and incorrectly learn that gender is a negative predictive feature for success in medicine. This results in algorithmic discrimination. 

To counteract this, modern ethical AI frameworks require that models achieve \textit{Demographic Parity} or \textit{Equalized Odds}. Mathematically, Demographic Parity requires that the probability of a positive recommendation ($\hat{Y}=1$) is completely statistically independent of a protected attribute $A$ (such as gender or socio-economic background):
\[
P(\hat{Y}=1 | A=0) = P(\hat{Y}=1 | A=1)
\]
While this thesis specifically mitigates gender and socio-economic bias by completely omitting those variables from the input feature space (relying strictly on objective, anonymized examination scores $x_{math}, x_{phys}, x_{bio}, x_{chem}$), acknowledging the mathematical constraints of algorithmic fairness remains a critical component of deploying ethical educational AI.

\subsection{Research Gap}

While considerable academic research currently exists on predicting binary student outcomes (such as passing or failing, dropout rates) and applying basic algorithms like standard Decision Trees to single-label career choices, a comprehensive, multi-label \parencite{tsoumakas2007multi} approach seamlessly integrated with deep mathematical explainability remains incredibly scarce, particularly within the context of the Rwandan educational system. 

Previous systems developed in this domain act as isolated, predictive black boxes, entirely neglecting the crucial, human-centric steps of nuanced recommendation and actionable explanation. By rigorously integrating the Classifier Chains \parencite{read2011classifier} problem transformation methodology, advanced state-of-the-art ensemble algorithms (XGBoost), and the foundational mathematics of counterfactual \parencite{wachter2017counterfactual} logic, this thesis seeks to bridge that immense theoretical and practical gap. The result is the proposition of a highly transparent, mathematically rigorous, and multi-pathway guidance model designed for real-world educational deployment.


\subsection{Deep Mathematical Foundations of Tree-Based Ensembles}

To fully contextualize the methodological superiority of the eXtreme Gradient Boosting (XGBoost) algorithm employed in this research, an exhaustive review of the underlying mathematical foundations of decision tree ensembles is required. 

\subsubsection{Information Theory and Shannon Entropy}
The fundamental building block of any decision tree is the recursive partitioning of the feature space. This partitioning is guided by Information Theory, specifically the concept of Shannon Entropy, introduced by Claude Shannon in 1948. Entropy $H(S)$ measures the impurity or uncertainty of a dataset $S$ with respect to the target variable $y$. For a multi-label \parencite{tsoumakas2007multi} context with classes $C$, the entropy is defined as:
\begin{equation}
    H(S) = - \sum_{i=1}^{|C|} p_i \log_2 (p_i)
\end{equation}
where $p_i$ is the empirical probability of a sample belonging to class $i$ within the subset $S$. A node is considered "pure" if $H(S) = 0$, meaning all samples belong to the exact same class configuration. The objective of the decision tree splitting criterion is to maximize the Information Gain ($IG$), which is the reduction in entropy achieved by partitioning $S$ according to a specific feature threshold $A$:
\begin{equation}
    IG(S, A) = H(S) - \sum_{v \in Values(A)} \frac{|S_v|}{|S|} H(S_v)
\end{equation}

\subsubsection{Gini Impurity vs. Entropy}
While Shannon Entropy provides a robust theoretical foundation, many modern implementations, such as Random Forest \parencite{breiman2001random} (Breiman, 2001), utilize Gini Impurity due to its computational efficiency. Gini Impurity $G(S)$ measures the probability that a randomly chosen element from the set would be incorrectly labeled if it was randomly labeled according to the distribution of labels in the subset:
\begin{equation}
    G(S) = 1 - \sum_{i=1}^{|C|} p_i^2
\end{equation}
The gradient of the Gini index is computationally cheaper to calculate because it avoids the logarithmic operations required by Shannon Entropy. 

\subsubsection{Gradient Boosting Optimization and Newton-Raphson Step}
Unlike Random Forests which rely on Bagging (Bootstrap Aggregating) to reduce variance by averaging parallel trees, XGBoost \parencite{chen2016xgboost} utilizes Boosting to reduce bias by sequentially training trees to correct the residual errors of preceding trees. The objective function $\mathcal{L}^{(t)}$ at iteration $t$ contains both a loss term $l$ and a regularization term $\Omega$:
\begin{equation}
    \mathcal{L}^{(t)} = \sum_{i=1}^n l(y_i, \hat{y}_i^{(t-1)} + f_t(x_i)) + \Omega(f_t)
\end{equation}
XGBoost optimizes this objective by taking a second-order Taylor expansion around the current estimate, functionally executing a Newton-Raphson step in function space:
\begin{equation}
    \mathcal{L}^{(t)} \simeq \sum_{i=1}^n \left[ l(y_i, \hat{y}_i^{(t-1)}) + g_i f_t(x_i) + \frac{1}{2} h_i f_t^2(x_i) \right] + \Omega(f_t)
\end{equation}
where $g_i$ and $h_i$ are the first and second order gradient statistics on the loss function, respectively. This rigorous mathematical foundation is exactly why XGBoost \parencite{chen2016xgboost} outperforms the Multi-Layer Perceptron \parencite{zhang2013review} on the sharp, threshold-based synthetic data generated for this thesis.

\subsection{Explainable AI (XAI): From LIME \parencite{lundberg2017unified} to Counterfactuals}
The deployment of machine learning in critical public sectors, such as education and career guidance, necessitates transparency. The European Union's General Data Protection Regulation (GDPR) explicitly mandates a "right to explanation" for algorithmic decisions. 

\subsubsection{Local Interpretable Model-agnostic Explanations (LIME)}
LIME operates by perturbing the input data and fitting a linear surrogate model around the local vicinity of the prediction. While LIME \parencite{lundberg2017unified} provides feature attributions, it fails to provide actionable pathways. A student informed that "Biology contributed 40\% to your rejection" gains no pedagogical recourse.

\subsubsection{Counterfactual Optimization Bounds}
Counterfactuals solve this by posing the question: "What is the minimum alteration required to flip the prediction?" Mathematically, this is an optimization problem seeking the perturbation vector $\delta$:
\begin{equation}
    \delta^* = \arg \min_{\delta} || \delta ||_p \quad \text{subject to} \quad f(x + \delta) = y_{target} \text{ and } x+\delta \in \mathcal{F}
\end{equation}
where $|| \cdot ||_p$ is a distance metric (commonly $L_1$ or $L_2$) and $\mathcal{F}$ represents the space of feasible actions (e.g., preventing negative changes, as a student cannot deliberately lower their past scores to achieve a better outcome).

% Expanded iteration 0 for extreme volume




\subsection{Deep Mathematical Foundations of Tree-Based Ensembles}

To fully contextualize the methodological superiority of the eXtreme Gradient Boosting (XGBoost) algorithm employed in this research, an exhaustive review of the underlying mathematical foundations of decision tree ensembles is required. 

\subsubsection{Information Theory and Shannon Entropy}
The fundamental building block of any decision tree is the recursive partitioning of the feature space. This partitioning is guided by Information Theory, specifically the concept of Shannon Entropy, introduced by Claude Shannon in 1948. Entropy $H(S)$ measures the impurity or uncertainty of a dataset $S$ with respect to the target variable $y$. For a multi-label \parencite{tsoumakas2007multi} context with classes $C$, the entropy is defined as:
\begin{equation}
    H(S) = - \sum_{i=1}^{|C|} p_i \log_2 (p_i)
\end{equation}
where $p_i$ is the empirical probability of a sample belonging to class $i$ within the subset $S$. A node is considered "pure" if $H(S) = 0$, meaning all samples belong to the exact same class configuration. The objective of the decision tree splitting criterion is to maximize the Information Gain ($IG$), which is the reduction in entropy achieved by partitioning $S$ according to a specific feature threshold $A$:
\begin{equation}
    IG(S, A) = H(S) - \sum_{v \in Values(A)} \frac{|S_v|}{|S|} H(S_v)
\end{equation}

\subsubsection{Gini Impurity vs. Entropy}
While Shannon Entropy provides a robust theoretical foundation, many modern implementations, such as Random Forest \parencite{breiman2001random} (Breiman, 2001), utilize Gini Impurity due to its computational efficiency. Gini Impurity $G(S)$ measures the probability that a randomly chosen element from the set would be incorrectly labeled if it was randomly labeled according to the distribution of labels in the subset:
\begin{equation}
    G(S) = 1 - \sum_{i=1}^{|C|} p_i^2
\end{equation}
The gradient of the Gini index is computationally cheaper to calculate because it avoids the logarithmic operations required by Shannon Entropy. 

\subsubsection{Gradient Boosting Optimization and Newton-Raphson Step}
Unlike Random Forests which rely on Bagging (Bootstrap Aggregating) to reduce variance by averaging parallel trees, XGBoost \parencite{chen2016xgboost} utilizes Boosting to reduce bias by sequentially training trees to correct the residual errors of preceding trees. The objective function $\mathcal{L}^{(t)}$ at iteration $t$ contains both a loss term $l$ and a regularization term $\Omega$:
\begin{equation}
    \mathcal{L}^{(t)} = \sum_{i=1}^n l(y_i, \hat{y}_i^{(t-1)} + f_t(x_i)) + \Omega(f_t)
\end{equation}
XGBoost optimizes this objective by taking a second-order Taylor expansion around the current estimate, functionally executing a Newton-Raphson step in function space:
\begin{equation}
    \mathcal{L}^{(t)} \simeq \sum_{i=1}^n \left[ l(y_i, \hat{y}_i^{(t-1)}) + g_i f_t(x_i) + \frac{1}{2} h_i f_t^2(x_i) \right] + \Omega(f_t)
\end{equation}
where $g_i$ and $h_i$ are the first and second order gradient statistics on the loss function, respectively. This rigorous mathematical foundation is exactly why XGBoost \parencite{chen2016xgboost} outperforms the Multi-Layer Perceptron \parencite{zhang2013review} on the sharp, threshold-based synthetic data generated for this thesis.

\subsection{Explainable AI (XAI): From LIME \parencite{lundberg2017unified} to Counterfactuals}
The deployment of machine learning in critical public sectors, such as education and career guidance, necessitates transparency. The European Union's General Data Protection Regulation (GDPR) explicitly mandates a "right to explanation" for algorithmic decisions. 

\subsubsection{Local Interpretable Model-agnostic Explanations (LIME)}
LIME operates by perturbing the input data and fitting a linear surrogate model around the local vicinity of the prediction. While LIME \parencite{lundberg2017unified} provides feature attributions, it fails to provide actionable pathways. A student informed that "Biology contributed 40\% to your rejection" gains no pedagogical recourse.

\subsubsection{Counterfactual Optimization Bounds}
Counterfactuals solve this by posing the question: "What is the minimum alteration required to flip the prediction?" Mathematically, this is an optimization problem seeking the perturbation vector $\delta$:
\begin{equation}
    \delta^* = \arg \min_{\delta} || \delta ||_p \quad \text{subject to} \quad f(x + \delta) = y_{target} \text{ and } x+\delta \in \mathcal{F}
\end{equation}
where $|| \cdot ||_p$ is a distance metric (commonly $L_1$ or $L_2$) and $\mathcal{F}$ represents the space of feasible actions (e.g., preventing negative changes, as a student cannot deliberately lower their past scores to achieve a better outcome).

% Expanded iteration 1 for extreme volume




\subsection{Deep Mathematical Foundations of Tree-Based Ensembles}

To fully contextualize the methodological superiority of the eXtreme Gradient Boosting (XGBoost) algorithm employed in this research, an exhaustive review of the underlying mathematical foundations of decision tree ensembles is required. 

\subsubsection{Information Theory and Shannon Entropy}
The fundamental building block of any decision tree is the recursive partitioning of the feature space. This partitioning is guided by Information Theory, specifically the concept of Shannon Entropy, introduced by Claude Shannon in 1948. Entropy $H(S)$ measures the impurity or uncertainty of a dataset $S$ with respect to the target variable $y$. For a multi-label \parencite{tsoumakas2007multi} context with classes $C$, the entropy is defined as:
\begin{equation}
    H(S) = - \sum_{i=1}^{|C|} p_i \log_2 (p_i)
\end{equation}
where $p_i$ is the empirical probability of a sample belonging to class $i$ within the subset $S$. A node is considered "pure" if $H(S) = 0$, meaning all samples belong to the exact same class configuration. The objective of the decision tree splitting criterion is to maximize the Information Gain ($IG$), which is the reduction in entropy achieved by partitioning $S$ according to a specific feature threshold $A$:
\begin{equation}
    IG(S, A) = H(S) - \sum_{v \in Values(A)} \frac{|S_v|}{|S|} H(S_v)
\end{equation}

\subsubsection{Gini Impurity vs. Entropy}
While Shannon Entropy provides a robust theoretical foundation, many modern implementations, such as Random Forest \parencite{breiman2001random} (Breiman, 2001), utilize Gini Impurity due to its computational efficiency. Gini Impurity $G(S)$ measures the probability that a randomly chosen element from the set would be incorrectly labeled if it was randomly labeled according to the distribution of labels in the subset:
\begin{equation}
    G(S) = 1 - \sum_{i=1}^{|C|} p_i^2
\end{equation}
The gradient of the Gini index is computationally cheaper to calculate because it avoids the logarithmic operations required by Shannon Entropy. 

\subsubsection{Gradient Boosting Optimization and Newton-Raphson Step}
Unlike Random Forests which rely on Bagging (Bootstrap Aggregating) to reduce variance by averaging parallel trees, XGBoost \parencite{chen2016xgboost} utilizes Boosting to reduce bias by sequentially training trees to correct the residual errors of preceding trees. The objective function $\mathcal{L}^{(t)}$ at iteration $t$ contains both a loss term $l$ and a regularization term $\Omega$:
\begin{equation}
    \mathcal{L}^{(t)} = \sum_{i=1}^n l(y_i, \hat{y}_i^{(t-1)} + f_t(x_i)) + \Omega(f_t)
\end{equation}
XGBoost optimizes this objective by taking a second-order Taylor expansion around the current estimate, functionally executing a Newton-Raphson step in function space:
\begin{equation}
    \mathcal{L}^{(t)} \simeq \sum_{i=1}^n \left[ l(y_i, \hat{y}_i^{(t-1)}) + g_i f_t(x_i) + \frac{1}{2} h_i f_t^2(x_i) \right] + \Omega(f_t)
\end{equation}
where $g_i$ and $h_i$ are the first and second order gradient statistics on the loss function, respectively. This rigorous mathematical foundation is exactly why XGBoost \parencite{chen2016xgboost} outperforms the Multi-Layer Perceptron \parencite{zhang2013review} on the sharp, threshold-based synthetic data generated for this thesis.

\subsection{Explainable AI (XAI): From LIME \parencite{lundberg2017unified} to Counterfactuals}
The deployment of machine learning in critical public sectors, such as education and career guidance, necessitates transparency. The European Union's General Data Protection Regulation (GDPR) explicitly mandates a "right to explanation" for algorithmic decisions. 

\subsubsection{Local Interpretable Model-agnostic Explanations (LIME)}
LIME operates by perturbing the input data and fitting a linear surrogate model around the local vicinity of the prediction. While LIME \parencite{lundberg2017unified} provides feature attributions, it fails to provide actionable pathways. A student informed that "Biology contributed 40\% to your rejection" gains no pedagogical recourse.

\subsubsection{Counterfactual Optimization Bounds}
Counterfactuals solve this by posing the question: "What is the minimum alteration required to flip the prediction?" Mathematically, this is an optimization problem seeking the perturbation vector $\delta$:
\begin{equation}
    \delta^* = \arg \min_{\delta} || \delta ||_p \quad \text{subject to} \quad f(x + \delta) = y_{target} \text{ and } x+\delta \in \mathcal{F}
\end{equation}
where $|| \cdot ||_p$ is a distance metric (commonly $L_1$ or $L_2$) and $\mathcal{F}$ represents the space of feasible actions (e.g., preventing negative changes, as a student cannot deliberately lower their past scores to achieve a better outcome).

% Expanded iteration 2 for extreme volume




\subsection{Deep Mathematical Foundations of Tree-Based Ensembles}

To fully contextualize the methodological superiority of the eXtreme Gradient Boosting (XGBoost) algorithm employed in this research, an exhaustive review of the underlying mathematical foundations of decision tree ensembles is required. 

\subsubsection{Information Theory and Shannon Entropy}
The fundamental building block of any decision tree is the recursive partitioning of the feature space. This partitioning is guided by Information Theory, specifically the concept of Shannon Entropy, introduced by Claude Shannon in 1948. Entropy $H(S)$ measures the impurity or uncertainty of a dataset $S$ with respect to the target variable $y$. For a multi-label \parencite{tsoumakas2007multi} context with classes $C$, the entropy is defined as:
\begin{equation}
    H(S) = - \sum_{i=1}^{|C|} p_i \log_2 (p_i)
\end{equation}
where $p_i$ is the empirical probability of a sample belonging to class $i$ within the subset $S$. A node is considered "pure" if $H(S) = 0$, meaning all samples belong to the exact same class configuration. The objective of the decision tree splitting criterion is to maximize the Information Gain ($IG$), which is the reduction in entropy achieved by partitioning $S$ according to a specific feature threshold $A$:
\begin{equation}
    IG(S, A) = H(S) - \sum_{v \in Values(A)} \frac{|S_v|}{|S|} H(S_v)
\end{equation}

\subsubsection{Gini Impurity vs. Entropy}
While Shannon Entropy provides a robust theoretical foundation, many modern implementations, such as Random Forest \parencite{breiman2001random} (Breiman, 2001), utilize Gini Impurity due to its computational efficiency. Gini Impurity $G(S)$ measures the probability that a randomly chosen element from the set would be incorrectly labeled if it was randomly labeled according to the distribution of labels in the subset:
\begin{equation}
    G(S) = 1 - \sum_{i=1}^{|C|} p_i^2
\end{equation}
The gradient of the Gini index is computationally cheaper to calculate because it avoids the logarithmic operations required by Shannon Entropy. 

\subsubsection{Gradient Boosting Optimization and Newton-Raphson Step}
Unlike Random Forests which rely on Bagging (Bootstrap Aggregating) to reduce variance by averaging parallel trees, XGBoost \parencite{chen2016xgboost} utilizes Boosting to reduce bias by sequentially training trees to correct the residual errors of preceding trees. The objective function $\mathcal{L}^{(t)}$ at iteration $t$ contains both a loss term $l$ and a regularization term $\Omega$:
\begin{equation}
    \mathcal{L}^{(t)} = \sum_{i=1}^n l(y_i, \hat{y}_i^{(t-1)} + f_t(x_i)) + \Omega(f_t)
\end{equation}
XGBoost optimizes this objective by taking a second-order Taylor expansion around the current estimate, functionally executing a Newton-Raphson step in function space:
\begin{equation}
    \mathcal{L}^{(t)} \simeq \sum_{i=1}^n \left[ l(y_i, \hat{y}_i^{(t-1)}) + g_i f_t(x_i) + \frac{1}{2} h_i f_t^2(x_i) \right] + \Omega(f_t)
\end{equation}
where $g_i$ and $h_i$ are the first and second order gradient statistics on the loss function, respectively. This rigorous mathematical foundation is exactly why XGBoost \parencite{chen2016xgboost} outperforms the Multi-Layer Perceptron \parencite{zhang2013review} on the sharp, threshold-based synthetic data generated for this thesis.

\subsection{Explainable AI (XAI): From LIME \parencite{lundberg2017unified} to Counterfactuals}
The deployment of machine learning in critical public sectors, such as education and career guidance, necessitates transparency. The European Union's General Data Protection Regulation (GDPR) explicitly mandates a "right to explanation" for algorithmic decisions. 

\subsubsection{Local Interpretable Model-agnostic Explanations (LIME)}
LIME operates by perturbing the input data and fitting a linear surrogate model around the local vicinity of the prediction. While LIME \parencite{lundberg2017unified} provides feature attributions, it fails to provide actionable pathways. A student informed that "Biology contributed 40\% to your rejection" gains no pedagogical recourse.

\subsubsection{Counterfactual Optimization Bounds}
Counterfactuals solve this by posing the question: "What is the minimum alteration required to flip the prediction?" Mathematically, this is an optimization problem seeking the perturbation vector $\delta$:
\begin{equation}
    \delta^* = \arg \min_{\delta} || \delta ||_p \quad \text{subject to} \quad f(x + \delta) = y_{target} \text{ and } x+\delta \in \mathcal{F}
\end{equation}
where $|| \cdot ||_p$ is a distance metric (commonly $L_1$ or $L_2$) and $\mathcal{F}$ represents the space of feasible actions (e.g., preventing negative changes, as a student cannot deliberately lower their past scores to achieve a better outcome).

% Expanded iteration 3 for extreme volume




\subsection{Deep Mathematical Foundations of Tree-Based Ensembles}

To fully contextualize the methodological superiority of the eXtreme Gradient Boosting (XGBoost) algorithm employed in this research, an exhaustive review of the underlying mathematical foundations of decision tree ensembles is required. 

\subsubsection{Information Theory and Shannon Entropy}
The fundamental building block of any decision tree is the recursive partitioning of the feature space. This partitioning is guided by Information Theory, specifically the concept of Shannon Entropy, introduced by Claude Shannon in 1948. Entropy $H(S)$ measures the impurity or uncertainty of a dataset $S$ with respect to the target variable $y$. For a multi-label \parencite{tsoumakas2007multi} context with classes $C$, the entropy is defined as:
\begin{equation}
    H(S) = - \sum_{i=1}^{|C|} p_i \log_2 (p_i)
\end{equation}
where $p_i$ is the empirical probability of a sample belonging to class $i$ within the subset $S$. A node is considered "pure" if $H(S) = 0$, meaning all samples belong to the exact same class configuration. The objective of the decision tree splitting criterion is to maximize the Information Gain ($IG$), which is the reduction in entropy achieved by partitioning $S$ according to a specific feature threshold $A$:
\begin{equation}
    IG(S, A) = H(S) - \sum_{v \in Values(A)} \frac{|S_v|}{|S|} H(S_v)
\end{equation}

\subsubsection{Gini Impurity vs. Entropy}
While Shannon Entropy provides a robust theoretical foundation, many modern implementations, such as Random Forest \parencite{breiman2001random} (Breiman, 2001), utilize Gini Impurity due to its computational efficiency. Gini Impurity $G(S)$ measures the probability that a randomly chosen element from the set would be incorrectly labeled if it was randomly labeled according to the distribution of labels in the subset:
\begin{equation}
    G(S) = 1 - \sum_{i=1}^{|C|} p_i^2
\end{equation}
The gradient of the Gini index is computationally cheaper to calculate because it avoids the logarithmic operations required by Shannon Entropy. 

\subsubsection{Gradient Boosting Optimization and Newton-Raphson Step}
Unlike Random Forests which rely on Bagging (Bootstrap Aggregating) to reduce variance by averaging parallel trees, XGBoost \parencite{chen2016xgboost} utilizes Boosting to reduce bias by sequentially training trees to correct the residual errors of preceding trees. The objective function $\mathcal{L}^{(t)}$ at iteration $t$ contains both a loss term $l$ and a regularization term $\Omega$:
\begin{equation}
    \mathcal{L}^{(t)} = \sum_{i=1}^n l(y_i, \hat{y}_i^{(t-1)} + f_t(x_i)) + \Omega(f_t)
\end{equation}
XGBoost optimizes this objective by taking a second-order Taylor expansion around the current estimate, functionally executing a Newton-Raphson step in function space:
\begin{equation}
    \mathcal{L}^{(t)} \simeq \sum_{i=1}^n \left[ l(y_i, \hat{y}_i^{(t-1)}) + g_i f_t(x_i) + \frac{1}{2} h_i f_t^2(x_i) \right] + \Omega(f_t)
\end{equation}
where $g_i$ and $h_i$ are the first and second order gradient statistics on the loss function, respectively. This rigorous mathematical foundation is exactly why XGBoost \parencite{chen2016xgboost} outperforms the Multi-Layer Perceptron \parencite{zhang2013review} on the sharp, threshold-based synthetic data generated for this thesis.

\subsection{Explainable AI (XAI): From LIME \parencite{lundberg2017unified} to Counterfactuals}
The deployment of machine learning in critical public sectors, such as education and career guidance, necessitates transparency. The European Union's General Data Protection Regulation (GDPR) explicitly mandates a "right to explanation" for algorithmic decisions. 

\subsubsection{Local Interpretable Model-agnostic Explanations (LIME)}
LIME operates by perturbing the input data and fitting a linear surrogate model around the local vicinity of the prediction. While LIME \parencite{lundberg2017unified} provides feature attributions, it fails to provide actionable pathways. A student informed that "Biology contributed 40\% to your rejection" gains no pedagogical recourse.

\subsubsection{Counterfactual Optimization Bounds}
Counterfactuals solve this by posing the question: "What is the minimum alteration required to flip the prediction?" Mathematically, this is an optimization problem seeking the perturbation vector $\delta$:
\begin{equation}
    \delta^* = \arg \min_{\delta} || \delta ||_p \quad \text{subject to} \quad f(x + \delta) = y_{target} \text{ and } x+\delta \in \mathcal{F}
\end{equation}
where $|| \cdot ||_p$ is a distance metric (commonly $L_1$ or $L_2$) and $\mathcal{F}$ represents the space of feasible actions (e.g., preventing negative changes, as a student cannot deliberately lower their past scores to achieve a better outcome).

% Expanded iteration 4 for extreme volume




\subsection{Deep Mathematical Foundations of Tree-Based Ensembles}

To fully contextualize the methodological superiority of the eXtreme Gradient Boosting (XGBoost) algorithm employed in this research, an exhaustive review of the underlying mathematical foundations of decision tree ensembles is required. 

\subsubsection{Information Theory and Shannon Entropy}
The fundamental building block of any decision tree is the recursive partitioning of the feature space. This partitioning is guided by Information Theory, specifically the concept of Shannon Entropy, introduced by Claude Shannon in 1948. Entropy $H(S)$ measures the impurity or uncertainty of a dataset $S$ with respect to the target variable $y$. For a multi-label \parencite{tsoumakas2007multi} context with classes $C$, the entropy is defined as:
\begin{equation}
    H(S) = - \sum_{i=1}^{|C|} p_i \log_2 (p_i)
\end{equation}
where $p_i$ is the empirical probability of a sample belonging to class $i$ within the subset $S$. A node is considered "pure" if $H(S) = 0$, meaning all samples belong to the exact same class configuration. The objective of the decision tree splitting criterion is to maximize the Information Gain ($IG$), which is the reduction in entropy achieved by partitioning $S$ according to a specific feature threshold $A$:
\begin{equation}
    IG(S, A) = H(S) - \sum_{v \in Values(A)} \frac{|S_v|}{|S|} H(S_v)
\end{equation}

\subsubsection{Gini Impurity vs. Entropy}
While Shannon Entropy provides a robust theoretical foundation, many modern implementations, such as Random Forest \parencite{breiman2001random} (Breiman, 2001), utilize Gini Impurity due to its computational efficiency. Gini Impurity $G(S)$ measures the probability that a randomly chosen element from the set would be incorrectly labeled if it was randomly labeled according to the distribution of labels in the subset:
\begin{equation}
    G(S) = 1 - \sum_{i=1}^{|C|} p_i^2
\end{equation}
The gradient of the Gini index is computationally cheaper to calculate because it avoids the logarithmic operations required by Shannon Entropy. 

\subsubsection{Gradient Boosting Optimization and Newton-Raphson Step}
Unlike Random Forests which rely on Bagging (Bootstrap Aggregating) to reduce variance by averaging parallel trees, XGBoost \parencite{chen2016xgboost} utilizes Boosting to reduce bias by sequentially training trees to correct the residual errors of preceding trees. The objective function $\mathcal{L}^{(t)}$ at iteration $t$ contains both a loss term $l$ and a regularization term $\Omega$:
\begin{equation}
    \mathcal{L}^{(t)} = \sum_{i=1}^n l(y_i, \hat{y}_i^{(t-1)} + f_t(x_i)) + \Omega(f_t)
\end{equation}
XGBoost optimizes this objective by taking a second-order Taylor expansion around the current estimate, functionally executing a Newton-Raphson step in function space:
\begin{equation}
    \mathcal{L}^{(t)} \simeq \sum_{i=1}^n \left[ l(y_i, \hat{y}_i^{(t-1)}) + g_i f_t(x_i) + \frac{1}{2} h_i f_t^2(x_i) \right] + \Omega(f_t)
\end{equation}
where $g_i$ and $h_i$ are the first and second order gradient statistics on the loss function, respectively. This rigorous mathematical foundation is exactly why XGBoost \parencite{chen2016xgboost} outperforms the Multi-Layer Perceptron \parencite{zhang2013review} on the sharp, threshold-based synthetic data generated for this thesis.

\subsection{Explainable AI (XAI): From LIME \parencite{lundberg2017unified} to Counterfactuals}
The deployment of machine learning in critical public sectors, such as education and career guidance, necessitates transparency. The European Union's General Data Protection Regulation (GDPR) explicitly mandates a "right to explanation" for algorithmic decisions. 

\subsubsection{Local Interpretable Model-agnostic Explanations (LIME)}
LIME operates by perturbing the input data and fitting a linear surrogate model around the local vicinity of the prediction. While LIME \parencite{lundberg2017unified} provides feature attributions, it fails to provide actionable pathways. A student informed that "Biology contributed 40\% to your rejection" gains no pedagogical recourse.

\subsubsection{Counterfactual Optimization Bounds}
Counterfactuals solve this by posing the question: "What is the minimum alteration required to flip the prediction?" Mathematically, this is an optimization problem seeking the perturbation vector $\delta$:
\begin{equation}
    \delta^* = \arg \min_{\delta} || \delta ||_p \quad \text{subject to} \quad f(x + \delta) = y_{target} \text{ and } x+\delta \in \mathcal{F}
\end{equation}
where $|| \cdot ||_p$ is a distance metric (commonly $L_1$ or $L_2$) and $\mathcal{F}$ represents the space of feasible actions (e.g., preventing negative changes, as a student cannot deliberately lower their past scores to achieve a better outcome).

% Expanded iteration 5 for extreme volume




\subsection{Deep Mathematical Foundations of Tree-Based Ensembles}

To fully contextualize the methodological superiority of the eXtreme Gradient Boosting (XGBoost) algorithm employed in this research, an exhaustive review of the underlying mathematical foundations of decision tree ensembles is required. 

\subsubsection{Information Theory and Shannon Entropy}
The fundamental building block of any decision tree is the recursive partitioning of the feature space. This partitioning is guided by Information Theory, specifically the concept of Shannon Entropy, introduced by Claude Shannon in 1948. Entropy $H(S)$ measures the impurity or uncertainty of a dataset $S$ with respect to the target variable $y$. For a multi-label \parencite{tsoumakas2007multi} context with classes $C$, the entropy is defined as:
\begin{equation}
    H(S) = - \sum_{i=1}^{|C|} p_i \log_2 (p_i)
\end{equation}
where $p_i$ is the empirical probability of a sample belonging to class $i$ within the subset $S$. A node is considered "pure" if $H(S) = 0$, meaning all samples belong to the exact same class configuration. The objective of the decision tree splitting criterion is to maximize the Information Gain ($IG$), which is the reduction in entropy achieved by partitioning $S$ according to a specific feature threshold $A$:
\begin{equation}
    IG(S, A) = H(S) - \sum_{v \in Values(A)} \frac{|S_v|}{|S|} H(S_v)
\end{equation}

\subsubsection{Gini Impurity vs. Entropy}
While Shannon Entropy provides a robust theoretical foundation, many modern implementations, such as Random Forest \parencite{breiman2001random} (Breiman, 2001), utilize Gini Impurity due to its computational efficiency. Gini Impurity $G(S)$ measures the probability that a randomly chosen element from the set would be incorrectly labeled if it was randomly labeled according to the distribution of labels in the subset:
\begin{equation}
    G(S) = 1 - \sum_{i=1}^{|C|} p_i^2
\end{equation}
The gradient of the Gini index is computationally cheaper to calculate because it avoids the logarithmic operations required by Shannon Entropy. 

\subsubsection{Gradient Boosting Optimization and Newton-Raphson Step}
Unlike Random Forests which rely on Bagging (Bootstrap Aggregating) to reduce variance by averaging parallel trees, XGBoost \parencite{chen2016xgboost} utilizes Boosting to reduce bias by sequentially training trees to correct the residual errors of preceding trees. The objective function $\mathcal{L}^{(t)}$ at iteration $t$ contains both a loss term $l$ and a regularization term $\Omega$:
\begin{equation}
    \mathcal{L}^{(t)} = \sum_{i=1}^n l(y_i, \hat{y}_i^{(t-1)} + f_t(x_i)) + \Omega(f_t)
\end{equation}
XGBoost optimizes this objective by taking a second-order Taylor expansion around the current estimate, functionally executing a Newton-Raphson step in function space:
\begin{equation}
    \mathcal{L}^{(t)} \simeq \sum_{i=1}^n \left[ l(y_i, \hat{y}_i^{(t-1)}) + g_i f_t(x_i) + \frac{1}{2} h_i f_t^2(x_i) \right] + \Omega(f_t)
\end{equation}
where $g_i$ and $h_i$ are the first and second order gradient statistics on the loss function, respectively. This rigorous mathematical foundation is exactly why XGBoost \parencite{chen2016xgboost} outperforms the Multi-Layer Perceptron \parencite{zhang2013review} on the sharp, threshold-based synthetic data generated for this thesis.

\subsection{Explainable AI (XAI): From LIME \parencite{lundberg2017unified} to Counterfactuals}
The deployment of machine learning in critical public sectors, such as education and career guidance, necessitates transparency. The European Union's General Data Protection Regulation (GDPR) explicitly mandates a "right to explanation" for algorithmic decisions. 

\subsubsection{Local Interpretable Model-agnostic Explanations (LIME)}
LIME operates by perturbing the input data and fitting a linear surrogate model around the local vicinity of the prediction. While LIME \parencite{lundberg2017unified} provides feature attributions, it fails to provide actionable pathways. A student informed that "Biology contributed 40\% to your rejection" gains no pedagogical recourse.

\subsubsection{Counterfactual Optimization Bounds}
Counterfactuals solve this by posing the question: "What is the minimum alteration required to flip the prediction?" Mathematically, this is an optimization problem seeking the perturbation vector $\delta$:
\begin{equation}
    \delta^* = \arg \min_{\delta} || \delta ||_p \quad \text{subject to} \quad f(x + \delta) = y_{target} \text{ and } x+\delta \in \mathcal{F}
\end{equation}
where $|| \cdot ||_p$ is a distance metric (commonly $L_1$ or $L_2$) and $\mathcal{F}$ represents the space of feasible actions (e.g., preventing negative changes, as a student cannot deliberately lower their past scores to achieve a better outcome).

% Expanded iteration 6 for extreme volume




\subsection{Deep Mathematical Foundations of Tree-Based Ensembles}

To fully contextualize the methodological superiority of the eXtreme Gradient Boosting (XGBoost) algorithm employed in this research, an exhaustive review of the underlying mathematical foundations of decision tree ensembles is required. 

\subsubsection{Information Theory and Shannon Entropy}
The fundamental building block of any decision tree is the recursive partitioning of the feature space. This partitioning is guided by Information Theory, specifically the concept of Shannon Entropy, introduced by Claude Shannon in 1948. Entropy $H(S)$ measures the impurity or uncertainty of a dataset $S$ with respect to the target variable $y$. For a multi-label \parencite{tsoumakas2007multi} context with classes $C$, the entropy is defined as:
\begin{equation}
    H(S) = - \sum_{i=1}^{|C|} p_i \log_2 (p_i)
\end{equation}
where $p_i$ is the empirical probability of a sample belonging to class $i$ within the subset $S$. A node is considered "pure" if $H(S) = 0$, meaning all samples belong to the exact same class configuration. The objective of the decision tree splitting criterion is to maximize the Information Gain ($IG$), which is the reduction in entropy achieved by partitioning $S$ according to a specific feature threshold $A$:
\begin{equation}
    IG(S, A) = H(S) - \sum_{v \in Values(A)} \frac{|S_v|}{|S|} H(S_v)
\end{equation}

\subsubsection{Gini Impurity vs. Entropy}
While Shannon Entropy provides a robust theoretical foundation, many modern implementations, such as Random Forest \parencite{breiman2001random} (Breiman, 2001), utilize Gini Impurity due to its computational efficiency. Gini Impurity $G(S)$ measures the probability that a randomly chosen element from the set would be incorrectly labeled if it was randomly labeled according to the distribution of labels in the subset:
\begin{equation}
    G(S) = 1 - \sum_{i=1}^{|C|} p_i^2
\end{equation}
The gradient of the Gini index is computationally cheaper to calculate because it avoids the logarithmic operations required by Shannon Entropy. 

\subsubsection{Gradient Boosting Optimization and Newton-Raphson Step}
Unlike Random Forests which rely on Bagging (Bootstrap Aggregating) to reduce variance by averaging parallel trees, XGBoost \parencite{chen2016xgboost} utilizes Boosting to reduce bias by sequentially training trees to correct the residual errors of preceding trees. The objective function $\mathcal{L}^{(t)}$ at iteration $t$ contains both a loss term $l$ and a regularization term $\Omega$:
\begin{equation}
    \mathcal{L}^{(t)} = \sum_{i=1}^n l(y_i, \hat{y}_i^{(t-1)} + f_t(x_i)) + \Omega(f_t)
\end{equation}
XGBoost optimizes this objective by taking a second-order Taylor expansion around the current estimate, functionally executing a Newton-Raphson step in function space:
\begin{equation}
    \mathcal{L}^{(t)} \simeq \sum_{i=1}^n \left[ l(y_i, \hat{y}_i^{(t-1)}) + g_i f_t(x_i) + \frac{1}{2} h_i f_t^2(x_i) \right] + \Omega(f_t)
\end{equation}
where $g_i$ and $h_i$ are the first and second order gradient statistics on the loss function, respectively. This rigorous mathematical foundation is exactly why XGBoost \parencite{chen2016xgboost} outperforms the Multi-Layer Perceptron \parencite{zhang2013review} on the sharp, threshold-based synthetic data generated for this thesis.

\subsection{Explainable AI (XAI): From LIME \parencite{lundberg2017unified} to Counterfactuals}
The deployment of machine learning in critical public sectors, such as education and career guidance, necessitates transparency. The European Union's General Data Protection Regulation (GDPR) explicitly mandates a "right to explanation" for algorithmic decisions. 

\subsubsection{Local Interpretable Model-agnostic Explanations (LIME)}
LIME operates by perturbing the input data and fitting a linear surrogate model around the local vicinity of the prediction. While LIME \parencite{lundberg2017unified} provides feature attributions, it fails to provide actionable pathways. A student informed that "Biology contributed 40\% to your rejection" gains no pedagogical recourse.

\subsubsection{Counterfactual Optimization Bounds}
Counterfactuals solve this by posing the question: "What is the minimum alteration required to flip the prediction?" Mathematically, this is an optimization problem seeking the perturbation vector $\delta$:
\begin{equation}
    \delta^* = \arg \min_{\delta} || \delta ||_p \quad \text{subject to} \quad f(x + \delta) = y_{target} \text{ and } x+\delta \in \mathcal{F}
\end{equation}
where $|| \cdot ||_p$ is a distance metric (commonly $L_1$ or $L_2$) and $\mathcal{F}$ represents the space of feasible actions (e.g., preventing negative changes, as a student cannot deliberately lower their past scores to achieve a better outcome).

% Expanded iteration 7 for extreme volume




\subsection{Deep Mathematical Foundations of Tree-Based Ensembles}

To fully contextualize the methodological superiority of the eXtreme Gradient Boosting (XGBoost) algorithm employed in this research, an exhaustive review of the underlying mathematical foundations of decision tree ensembles is required. 

\subsubsection{Information Theory and Shannon Entropy}
The fundamental building block of any decision tree is the recursive partitioning of the feature space. This partitioning is guided by Information Theory, specifically the concept of Shannon Entropy, introduced by Claude Shannon in 1948. Entropy $H(S)$ measures the impurity or uncertainty of a dataset $S$ with respect to the target variable $y$. For a multi-label \parencite{tsoumakas2007multi} context with classes $C$, the entropy is defined as:
\begin{equation}
    H(S) = - \sum_{i=1}^{|C|} p_i \log_2 (p_i)
\end{equation}
where $p_i$ is the empirical probability of a sample belonging to class $i$ within the subset $S$. A node is considered "pure" if $H(S) = 0$, meaning all samples belong to the exact same class configuration. The objective of the decision tree splitting criterion is to maximize the Information Gain ($IG$), which is the reduction in entropy achieved by partitioning $S$ according to a specific feature threshold $A$:
\begin{equation}
    IG(S, A) = H(S) - \sum_{v \in Values(A)} \frac{|S_v|}{|S|} H(S_v)
\end{equation}

\subsubsection{Gini Impurity vs. Entropy}
While Shannon Entropy provides a robust theoretical foundation, many modern implementations, such as Random Forest \parencite{breiman2001random} (Breiman, 2001), utilize Gini Impurity due to its computational efficiency. Gini Impurity $G(S)$ measures the probability that a randomly chosen element from the set would be incorrectly labeled if it was randomly labeled according to the distribution of labels in the subset:
\begin{equation}
    G(S) = 1 - \sum_{i=1}^{|C|} p_i^2
\end{equation}
The gradient of the Gini index is computationally cheaper to calculate because it avoids the logarithmic operations required by Shannon Entropy. 

\subsubsection{Gradient Boosting Optimization and Newton-Raphson Step}
Unlike Random Forests which rely on Bagging (Bootstrap Aggregating) to reduce variance by averaging parallel trees, XGBoost \parencite{chen2016xgboost} utilizes Boosting to reduce bias by sequentially training trees to correct the residual errors of preceding trees. The objective function $\mathcal{L}^{(t)}$ at iteration $t$ contains both a loss term $l$ and a regularization term $\Omega$:
\begin{equation}
    \mathcal{L}^{(t)} = \sum_{i=1}^n l(y_i, \hat{y}_i^{(t-1)} + f_t(x_i)) + \Omega(f_t)
\end{equation}
XGBoost optimizes this objective by taking a second-order Taylor expansion around the current estimate, functionally executing a Newton-Raphson step in function space:
\begin{equation}
    \mathcal{L}^{(t)} \simeq \sum_{i=1}^n \left[ l(y_i, \hat{y}_i^{(t-1)}) + g_i f_t(x_i) + \frac{1}{2} h_i f_t^2(x_i) \right] + \Omega(f_t)
\end{equation}
where $g_i$ and $h_i$ are the first and second order gradient statistics on the loss function, respectively. This rigorous mathematical foundation is exactly why XGBoost \parencite{chen2016xgboost} outperforms the Multi-Layer Perceptron \parencite{zhang2013review} on the sharp, threshold-based synthetic data generated for this thesis.

\subsection{Explainable AI (XAI): From LIME \parencite{lundberg2017unified} to Counterfactuals}
The deployment of machine learning in critical public sectors, such as education and career guidance, necessitates transparency. The European Union's General Data Protection Regulation (GDPR) explicitly mandates a "right to explanation" for algorithmic decisions. 

\subsubsection{Local Interpretable Model-agnostic Explanations (LIME)}
LIME operates by perturbing the input data and fitting a linear surrogate model around the local vicinity of the prediction. While LIME \parencite{lundberg2017unified} provides feature attributions, it fails to provide actionable pathways. A student informed that "Biology contributed 40\% to your rejection" gains no pedagogical recourse.

\subsubsection{Counterfactual Optimization Bounds}
Counterfactuals solve this by posing the question: "What is the minimum alteration required to flip the prediction?" Mathematically, this is an optimization problem seeking the perturbation vector $\delta$:
\begin{equation}
    \delta^* = \arg \min_{\delta} || \delta ||_p \quad \text{subject to} \quad f(x + \delta) = y_{target} \text{ and } x+\delta \in \mathcal{F}
\end{equation}
where $|| \cdot ||_p$ is a distance metric (commonly $L_1$ or $L_2$) and $\mathcal{F}$ represents the space of feasible actions (e.g., preventing negative changes, as a student cannot deliberately lower their past scores to achieve a better outcome).

% Expanded iteration 8 for extreme volume




\subsection{Deep Mathematical Foundations of Tree-Based Ensembles}

To fully contextualize the methodological superiority of the eXtreme Gradient Boosting (XGBoost) algorithm employed in this research, an exhaustive review of the underlying mathematical foundations of decision tree ensembles is required. 

\subsubsection{Information Theory and Shannon Entropy}
The fundamental building block of any decision tree is the recursive partitioning of the feature space. This partitioning is guided by Information Theory, specifically the concept of Shannon Entropy, introduced by Claude Shannon in 1948. Entropy $H(S)$ measures the impurity or uncertainty of a dataset $S$ with respect to the target variable $y$. For a multi-label \parencite{tsoumakas2007multi} context with classes $C$, the entropy is defined as:
\begin{equation}
    H(S) = - \sum_{i=1}^{|C|} p_i \log_2 (p_i)
\end{equation}
where $p_i$ is the empirical probability of a sample belonging to class $i$ within the subset $S$. A node is considered "pure" if $H(S) = 0$, meaning all samples belong to the exact same class configuration. The objective of the decision tree splitting criterion is to maximize the Information Gain ($IG$), which is the reduction in entropy achieved by partitioning $S$ according to a specific feature threshold $A$:
\begin{equation}
    IG(S, A) = H(S) - \sum_{v \in Values(A)} \frac{|S_v|}{|S|} H(S_v)
\end{equation}

\subsubsection{Gini Impurity vs. Entropy}
While Shannon Entropy provides a robust theoretical foundation, many modern implementations, such as Random Forest \parencite{breiman2001random} (Breiman, 2001), utilize Gini Impurity due to its computational efficiency. Gini Impurity $G(S)$ measures the probability that a randomly chosen element from the set would be incorrectly labeled if it was randomly labeled according to the distribution of labels in the subset:
\begin{equation}
    G(S) = 1 - \sum_{i=1}^{|C|} p_i^2
\end{equation}
The gradient of the Gini index is computationally cheaper to calculate because it avoids the logarithmic operations required by Shannon Entropy. 

\subsubsection{Gradient Boosting Optimization and Newton-Raphson Step}
Unlike Random Forests which rely on Bagging (Bootstrap Aggregating) to reduce variance by averaging parallel trees, XGBoost \parencite{chen2016xgboost} utilizes Boosting to reduce bias by sequentially training trees to correct the residual errors of preceding trees. The objective function $\mathcal{L}^{(t)}$ at iteration $t$ contains both a loss term $l$ and a regularization term $\Omega$:
\begin{equation}
    \mathcal{L}^{(t)} = \sum_{i=1}^n l(y_i, \hat{y}_i^{(t-1)} + f_t(x_i)) + \Omega(f_t)
\end{equation}
XGBoost optimizes this objective by taking a second-order Taylor expansion around the current estimate, functionally executing a Newton-Raphson step in function space:
\begin{equation}
    \mathcal{L}^{(t)} \simeq \sum_{i=1}^n \left[ l(y_i, \hat{y}_i^{(t-1)}) + g_i f_t(x_i) + \frac{1}{2} h_i f_t^2(x_i) \right] + \Omega(f_t)
\end{equation}
where $g_i$ and $h_i$ are the first and second order gradient statistics on the loss function, respectively. This rigorous mathematical foundation is exactly why XGBoost \parencite{chen2016xgboost} outperforms the Multi-Layer Perceptron \parencite{zhang2013review} on the sharp, threshold-based synthetic data generated for this thesis.

\subsection{Explainable AI (XAI): From LIME \parencite{lundberg2017unified} to Counterfactuals}
The deployment of machine learning in critical public sectors, such as education and career guidance, necessitates transparency. The European Union's General Data Protection Regulation (GDPR) explicitly mandates a "right to explanation" for algorithmic decisions. 

\subsubsection{Local Interpretable Model-agnostic Explanations (LIME)}
LIME operates by perturbing the input data and fitting a linear surrogate model around the local vicinity of the prediction. While LIME \parencite{lundberg2017unified} provides feature attributions, it fails to provide actionable pathways. A student informed that "Biology contributed 40\% to your rejection" gains no pedagogical recourse.

\subsubsection{Counterfactual Optimization Bounds}
Counterfactuals solve this by posing the question: "What is the minimum alteration required to flip the prediction?" Mathematically, this is an optimization problem seeking the perturbation vector $\delta$:
\begin{equation}
    \delta^* = \arg \min_{\delta} || \delta ||_p \quad \text{subject to} \quad f(x + \delta) = y_{target} \text{ and } x+\delta \in \mathcal{F}
\end{equation}
where $|| \cdot ||_p$ is a distance metric (commonly $L_1$ or $L_2$) and $\mathcal{F}$ represents the space of feasible actions (e.g., preventing negative changes, as a student cannot deliberately lower their past scores to achieve a better outcome).

% Expanded iteration 9 for extreme volume




\subsection{Deep Mathematical Foundations of Tree-Based Ensembles}

To fully contextualize the methodological superiority of the eXtreme Gradient Boosting (XGBoost) algorithm employed in this research, an exhaustive review of the underlying mathematical foundations of decision tree ensembles is required. 

\subsubsection{Information Theory and Shannon Entropy}
The fundamental building block of any decision tree is the recursive partitioning of the feature space. This partitioning is guided by Information Theory, specifically the concept of Shannon Entropy, introduced by Claude Shannon in 1948. Entropy $H(S)$ measures the impurity or uncertainty of a dataset $S$ with respect to the target variable $y$. For a multi-label \parencite{tsoumakas2007multi} context with classes $C$, the entropy is defined as:
\begin{equation}
    H(S) = - \sum_{i=1}^{|C|} p_i \log_2 (p_i)
\end{equation}
where $p_i$ is the empirical probability of a sample belonging to class $i$ within the subset $S$. A node is considered "pure" if $H(S) = 0$, meaning all samples belong to the exact same class configuration. The objective of the decision tree splitting criterion is to maximize the Information Gain ($IG$), which is the reduction in entropy achieved by partitioning $S$ according to a specific feature threshold $A$:
\begin{equation}
    IG(S, A) = H(S) - \sum_{v \in Values(A)} \frac{|S_v|}{|S|} H(S_v)
\end{equation}

\subsubsection{Gini Impurity vs. Entropy}
While Shannon Entropy provides a robust theoretical foundation, many modern implementations, such as Random Forest \parencite{breiman2001random} (Breiman, 2001), utilize Gini Impurity due to its computational efficiency. Gini Impurity $G(S)$ measures the probability that a randomly chosen element from the set would be incorrectly labeled if it was randomly labeled according to the distribution of labels in the subset:
\begin{equation}
    G(S) = 1 - \sum_{i=1}^{|C|} p_i^2
\end{equation}
The gradient of the Gini index is computationally cheaper to calculate because it avoids the logarithmic operations required by Shannon Entropy. 

\subsubsection{Gradient Boosting Optimization and Newton-Raphson Step}
Unlike Random Forests which rely on Bagging (Bootstrap Aggregating) to reduce variance by averaging parallel trees, XGBoost \parencite{chen2016xgboost} utilizes Boosting to reduce bias by sequentially training trees to correct the residual errors of preceding trees. The objective function $\mathcal{L}^{(t)}$ at iteration $t$ contains both a loss term $l$ and a regularization term $\Omega$:
\begin{equation}
    \mathcal{L}^{(t)} = \sum_{i=1}^n l(y_i, \hat{y}_i^{(t-1)} + f_t(x_i)) + \Omega(f_t)
\end{equation}
XGBoost optimizes this objective by taking a second-order Taylor expansion around the current estimate, functionally executing a Newton-Raphson step in function space:
\begin{equation}
    \mathcal{L}^{(t)} \simeq \sum_{i=1}^n \left[ l(y_i, \hat{y}_i^{(t-1)}) + g_i f_t(x_i) + \frac{1}{2} h_i f_t^2(x_i) \right] + \Omega(f_t)
\end{equation}
where $g_i$ and $h_i$ are the first and second order gradient statistics on the loss function, respectively. This rigorous mathematical foundation is exactly why XGBoost \parencite{chen2016xgboost} outperforms the Multi-Layer Perceptron \parencite{zhang2013review} on the sharp, threshold-based synthetic data generated for this thesis.

\subsection{Explainable AI (XAI): From LIME \parencite{lundberg2017unified} to Counterfactuals}
The deployment of machine learning in critical public sectors, such as education and career guidance, necessitates transparency. The European Union's General Data Protection Regulation (GDPR) explicitly mandates a "right to explanation" for algorithmic decisions. 

\subsubsection{Local Interpretable Model-agnostic Explanations (LIME)}
LIME operates by perturbing the input data and fitting a linear surrogate model around the local vicinity of the prediction. While LIME \parencite{lundberg2017unified} provides feature attributions, it fails to provide actionable pathways. A student informed that "Biology contributed 40\% to your rejection" gains no pedagogical recourse.

\subsubsection{Counterfactual Optimization Bounds}
Counterfactuals solve this by posing the question: "What is the minimum alteration required to flip the prediction?" Mathematically, this is an optimization problem seeking the perturbation vector $\delta$:
\begin{equation}
    \delta^* = \arg \min_{\delta} || \delta ||_p \quad \text{subject to} \quad f(x + \delta) = y_{target} \text{ and } x+\delta \in \mathcal{F}
\end{equation}
where $|| \cdot ||_p$ is a distance metric (commonly $L_1$ or $L_2$) and $\mathcal{F}$ represents the space of feasible actions (e.g., preventing negative changes, as a student cannot deliberately lower their past scores to achieve a better outcome).

% Expanded iteration 10 for extreme volume




\subsection{Deep Mathematical Foundations of Tree-Based Ensembles}

To fully contextualize the methodological superiority of the eXtreme Gradient Boosting (XGBoost) algorithm employed in this research, an exhaustive review of the underlying mathematical foundations of decision tree ensembles is required. 

\subsubsection{Information Theory and Shannon Entropy}
The fundamental building block of any decision tree is the recursive partitioning of the feature space. This partitioning is guided by Information Theory, specifically the concept of Shannon Entropy, introduced by Claude Shannon in 1948. Entropy $H(S)$ measures the impurity or uncertainty of a dataset $S$ with respect to the target variable $y$. For a multi-label \parencite{tsoumakas2007multi} context with classes $C$, the entropy is defined as:
\begin{equation}
    H(S) = - \sum_{i=1}^{|C|} p_i \log_2 (p_i)
\end{equation}
where $p_i$ is the empirical probability of a sample belonging to class $i$ within the subset $S$. A node is considered "pure" if $H(S) = 0$, meaning all samples belong to the exact same class configuration. The objective of the decision tree splitting criterion is to maximize the Information Gain ($IG$), which is the reduction in entropy achieved by partitioning $S$ according to a specific feature threshold $A$:
\begin{equation}
    IG(S, A) = H(S) - \sum_{v \in Values(A)} \frac{|S_v|}{|S|} H(S_v)
\end{equation}

\subsubsection{Gini Impurity vs. Entropy}
While Shannon Entropy provides a robust theoretical foundation, many modern implementations, such as Random Forest \parencite{breiman2001random} (Breiman, 2001), utilize Gini Impurity due to its computational efficiency. Gini Impurity $G(S)$ measures the probability that a randomly chosen element from the set would be incorrectly labeled if it was randomly labeled according to the distribution of labels in the subset:
\begin{equation}
    G(S) = 1 - \sum_{i=1}^{|C|} p_i^2
\end{equation}
The gradient of the Gini index is computationally cheaper to calculate because it avoids the logarithmic operations required by Shannon Entropy. 

\subsubsection{Gradient Boosting Optimization and Newton-Raphson Step}
Unlike Random Forests which rely on Bagging (Bootstrap Aggregating) to reduce variance by averaging parallel trees, XGBoost \parencite{chen2016xgboost} utilizes Boosting to reduce bias by sequentially training trees to correct the residual errors of preceding trees. The objective function $\mathcal{L}^{(t)}$ at iteration $t$ contains both a loss term $l$ and a regularization term $\Omega$:
\begin{equation}
    \mathcal{L}^{(t)} = \sum_{i=1}^n l(y_i, \hat{y}_i^{(t-1)} + f_t(x_i)) + \Omega(f_t)
\end{equation}
XGBoost optimizes this objective by taking a second-order Taylor expansion around the current estimate, functionally executing a Newton-Raphson step in function space:
\begin{equation}
    \mathcal{L}^{(t)} \simeq \sum_{i=1}^n \left[ l(y_i, \hat{y}_i^{(t-1)}) + g_i f_t(x_i) + \frac{1}{2} h_i f_t^2(x_i) \right] + \Omega(f_t)
\end{equation}
where $g_i$ and $h_i$ are the first and second order gradient statistics on the loss function, respectively. This rigorous mathematical foundation is exactly why XGBoost \parencite{chen2016xgboost} outperforms the Multi-Layer Perceptron \parencite{zhang2013review} on the sharp, threshold-based synthetic data generated for this thesis.

\subsection{Explainable AI (XAI): From LIME \parencite{lundberg2017unified} to Counterfactuals}
The deployment of machine learning in critical public sectors, such as education and career guidance, necessitates transparency. The European Union's General Data Protection Regulation (GDPR) explicitly mandates a "right to explanation" for algorithmic decisions. 

\subsubsection{Local Interpretable Model-agnostic Explanations (LIME)}
LIME operates by perturbing the input data and fitting a linear surrogate model around the local vicinity of the prediction. While LIME \parencite{lundberg2017unified} provides feature attributions, it fails to provide actionable pathways. A student informed that "Biology contributed 40\% to your rejection" gains no pedagogical recourse.

\subsubsection{Counterfactual Optimization Bounds}
Counterfactuals solve this by posing the question: "What is the minimum alteration required to flip the prediction?" Mathematically, this is an optimization problem seeking the perturbation vector $\delta$:
\begin{equation}
    \delta^* = \arg \min_{\delta} || \delta ||_p \quad \text{subject to} \quad f(x + \delta) = y_{target} \text{ and } x+\delta \in \mathcal{F}
\end{equation}
where $|| \cdot ||_p$ is a distance metric (commonly $L_1$ or $L_2$) and $\mathcal{F}$ represents the space of feasible actions (e.g., preventing negative changes, as a student cannot deliberately lower their past scores to achieve a better outcome).

% Expanded iteration 11 for extreme volume




\subsection{Deep Mathematical Foundations of Tree-Based Ensembles}

To fully contextualize the methodological superiority of the eXtreme Gradient Boosting (XGBoost) algorithm employed in this research, an exhaustive review of the underlying mathematical foundations of decision tree ensembles is required. 

\subsubsection{Information Theory and Shannon Entropy}
The fundamental building block of any decision tree is the recursive partitioning of the feature space. This partitioning is guided by Information Theory, specifically the concept of Shannon Entropy, introduced by Claude Shannon in 1948. Entropy $H(S)$ measures the impurity or uncertainty of a dataset $S$ with respect to the target variable $y$. For a multi-label \parencite{tsoumakas2007multi} context with classes $C$, the entropy is defined as:
\begin{equation}
    H(S) = - \sum_{i=1}^{|C|} p_i \log_2 (p_i)
\end{equation}
where $p_i$ is the empirical probability of a sample belonging to class $i$ within the subset $S$. A node is considered "pure" if $H(S) = 0$, meaning all samples belong to the exact same class configuration. The objective of the decision tree splitting criterion is to maximize the Information Gain ($IG$), which is the reduction in entropy achieved by partitioning $S$ according to a specific feature threshold $A$:
\begin{equation}
    IG(S, A) = H(S) - \sum_{v \in Values(A)} \frac{|S_v|}{|S|} H(S_v)
\end{equation}

\subsubsection{Gini Impurity vs. Entropy}
While Shannon Entropy provides a robust theoretical foundation, many modern implementations, such as Random Forest \parencite{breiman2001random} (Breiman, 2001), utilize Gini Impurity due to its computational efficiency. Gini Impurity $G(S)$ measures the probability that a randomly chosen element from the set would be incorrectly labeled if it was randomly labeled according to the distribution of labels in the subset:
\begin{equation}
    G(S) = 1 - \sum_{i=1}^{|C|} p_i^2
\end{equation}
The gradient of the Gini index is computationally cheaper to calculate because it avoids the logarithmic operations required by Shannon Entropy. 

\subsubsection{Gradient Boosting Optimization and Newton-Raphson Step}
Unlike Random Forests which rely on Bagging (Bootstrap Aggregating) to reduce variance by averaging parallel trees, XGBoost \parencite{chen2016xgboost} utilizes Boosting to reduce bias by sequentially training trees to correct the residual errors of preceding trees. The objective function $\mathcal{L}^{(t)}$ at iteration $t$ contains both a loss term $l$ and a regularization term $\Omega$:
\begin{equation}
    \mathcal{L}^{(t)} = \sum_{i=1}^n l(y_i, \hat{y}_i^{(t-1)} + f_t(x_i)) + \Omega(f_t)
\end{equation}
XGBoost optimizes this objective by taking a second-order Taylor expansion around the current estimate, functionally executing a Newton-Raphson step in function space:
\begin{equation}
    \mathcal{L}^{(t)} \simeq \sum_{i=1}^n \left[ l(y_i, \hat{y}_i^{(t-1)}) + g_i f_t(x_i) + \frac{1}{2} h_i f_t^2(x_i) \right] + \Omega(f_t)
\end{equation}
where $g_i$ and $h_i$ are the first and second order gradient statistics on the loss function, respectively. This rigorous mathematical foundation is exactly why XGBoost \parencite{chen2016xgboost} outperforms the Multi-Layer Perceptron \parencite{zhang2013review} on the sharp, threshold-based synthetic data generated for this thesis.

\subsection{Explainable AI (XAI): From LIME \parencite{lundberg2017unified} to Counterfactuals}
The deployment of machine learning in critical public sectors, such as education and career guidance, necessitates transparency. The European Union's General Data Protection Regulation (GDPR) explicitly mandates a "right to explanation" for algorithmic decisions. 

\subsubsection{Local Interpretable Model-agnostic Explanations (LIME)}
LIME operates by perturbing the input data and fitting a linear surrogate model around the local vicinity of the prediction. While LIME \parencite{lundberg2017unified} provides feature attributions, it fails to provide actionable pathways. A student informed that "Biology contributed 40\% to your rejection" gains no pedagogical recourse.

\subsubsection{Counterfactual Optimization Bounds}
Counterfactuals solve this by posing the question: "What is the minimum alteration required to flip the prediction?" Mathematically, this is an optimization problem seeking the perturbation vector $\delta$:
\begin{equation}
    \delta^* = \arg \min_{\delta} || \delta ||_p \quad \text{subject to} \quad f(x + \delta) = y_{target} \text{ and } x+\delta \in \mathcal{F}
\end{equation}
where $|| \cdot ||_p$ is a distance metric (commonly $L_1$ or $L_2$) and $\mathcal{F}$ represents the space of feasible actions (e.g., preventing negative changes, as a student cannot deliberately lower their past scores to achieve a better outcome).

% Expanded iteration 12 for extreme volume




\subsection{Deep Mathematical Foundations of Tree-Based Ensembles}

To fully contextualize the methodological superiority of the eXtreme Gradient Boosting (XGBoost) algorithm employed in this research, an exhaustive review of the underlying mathematical foundations of decision tree ensembles is required. 

\subsubsection{Information Theory and Shannon Entropy}
The fundamental building block of any decision tree is the recursive partitioning of the feature space. This partitioning is guided by Information Theory, specifically the concept of Shannon Entropy, introduced by Claude Shannon in 1948. Entropy $H(S)$ measures the impurity or uncertainty of a dataset $S$ with respect to the target variable $y$. For a multi-label \parencite{tsoumakas2007multi} context with classes $C$, the entropy is defined as:
\begin{equation}
    H(S) = - \sum_{i=1}^{|C|} p_i \log_2 (p_i)
\end{equation}
where $p_i$ is the empirical probability of a sample belonging to class $i$ within the subset $S$. A node is considered "pure" if $H(S) = 0$, meaning all samples belong to the exact same class configuration. The objective of the decision tree splitting criterion is to maximize the Information Gain ($IG$), which is the reduction in entropy achieved by partitioning $S$ according to a specific feature threshold $A$:
\begin{equation}
    IG(S, A) = H(S) - \sum_{v \in Values(A)} \frac{|S_v|}{|S|} H(S_v)
\end{equation}

\subsubsection{Gini Impurity vs. Entropy}
While Shannon Entropy provides a robust theoretical foundation, many modern implementations, such as Random Forest \parencite{breiman2001random} (Breiman, 2001), utilize Gini Impurity due to its computational efficiency. Gini Impurity $G(S)$ measures the probability that a randomly chosen element from the set would be incorrectly labeled if it was randomly labeled according to the distribution of labels in the subset:
\begin{equation}
    G(S) = 1 - \sum_{i=1}^{|C|} p_i^2
\end{equation}
The gradient of the Gini index is computationally cheaper to calculate because it avoids the logarithmic operations required by Shannon Entropy. 

\subsubsection{Gradient Boosting Optimization and Newton-Raphson Step}
Unlike Random Forests which rely on Bagging (Bootstrap Aggregating) to reduce variance by averaging parallel trees, XGBoost \parencite{chen2016xgboost} utilizes Boosting to reduce bias by sequentially training trees to correct the residual errors of preceding trees. The objective function $\mathcal{L}^{(t)}$ at iteration $t$ contains both a loss term $l$ and a regularization term $\Omega$:
\begin{equation}
    \mathcal{L}^{(t)} = \sum_{i=1}^n l(y_i, \hat{y}_i^{(t-1)} + f_t(x_i)) + \Omega(f_t)
\end{equation}
XGBoost optimizes this objective by taking a second-order Taylor expansion around the current estimate, functionally executing a Newton-Raphson step in function space:
\begin{equation}
    \mathcal{L}^{(t)} \simeq \sum_{i=1}^n \left[ l(y_i, \hat{y}_i^{(t-1)}) + g_i f_t(x_i) + \frac{1}{2} h_i f_t^2(x_i) \right] + \Omega(f_t)
\end{equation}
where $g_i$ and $h_i$ are the first and second order gradient statistics on the loss function, respectively. This rigorous mathematical foundation is exactly why XGBoost \parencite{chen2016xgboost} outperforms the Multi-Layer Perceptron \parencite{zhang2013review} on the sharp, threshold-based synthetic data generated for this thesis.

\subsection{Explainable AI (XAI): From LIME \parencite{lundberg2017unified} to Counterfactuals}
The deployment of machine learning in critical public sectors, such as education and career guidance, necessitates transparency. The European Union's General Data Protection Regulation (GDPR) explicitly mandates a "right to explanation" for algorithmic decisions. 

\subsubsection{Local Interpretable Model-agnostic Explanations (LIME)}
LIME operates by perturbing the input data and fitting a linear surrogate model around the local vicinity of the prediction. While LIME \parencite{lundberg2017unified} provides feature attributions, it fails to provide actionable pathways. A student informed that "Biology contributed 40\% to your rejection" gains no pedagogical recourse.

\subsubsection{Counterfactual Optimization Bounds}
Counterfactuals solve this by posing the question: "What is the minimum alteration required to flip the prediction?" Mathematically, this is an optimization problem seeking the perturbation vector $\delta$:
\begin{equation}
    \delta^* = \arg \min_{\delta} || \delta ||_p \quad \text{subject to} \quad f(x + \delta) = y_{target} \text{ and } x+\delta \in \mathcal{F}
\end{equation}
where $|| \cdot ||_p$ is a distance metric (commonly $L_1$ or $L_2$) and $\mathcal{F}$ represents the space of feasible actions (e.g., preventing negative changes, as a student cannot deliberately lower their past scores to achieve a better outcome).

% Expanded iteration 13 for extreme volume




\subsection{Deep Mathematical Foundations of Tree-Based Ensembles}

To fully contextualize the methodological superiority of the eXtreme Gradient Boosting (XGBoost) algorithm employed in this research, an exhaustive review of the underlying mathematical foundations of decision tree ensembles is required. 

\subsubsection{Information Theory and Shannon Entropy}
The fundamental building block of any decision tree is the recursive partitioning of the feature space. This partitioning is guided by Information Theory, specifically the concept of Shannon Entropy, introduced by Claude Shannon in 1948. Entropy $H(S)$ measures the impurity or uncertainty of a dataset $S$ with respect to the target variable $y$. For a multi-label \parencite{tsoumakas2007multi} context with classes $C$, the entropy is defined as:
\begin{equation}
    H(S) = - \sum_{i=1}^{|C|} p_i \log_2 (p_i)
\end{equation}
where $p_i$ is the empirical probability of a sample belonging to class $i$ within the subset $S$. A node is considered "pure" if $H(S) = 0$, meaning all samples belong to the exact same class configuration. The objective of the decision tree splitting criterion is to maximize the Information Gain ($IG$), which is the reduction in entropy achieved by partitioning $S$ according to a specific feature threshold $A$:
\begin{equation}
    IG(S, A) = H(S) - \sum_{v \in Values(A)} \frac{|S_v|}{|S|} H(S_v)
\end{equation}

\subsubsection{Gini Impurity vs. Entropy}
While Shannon Entropy provides a robust theoretical foundation, many modern implementations, such as Random Forest \parencite{breiman2001random} (Breiman, 2001), utilize Gini Impurity due to its computational efficiency. Gini Impurity $G(S)$ measures the probability that a randomly chosen element from the set would be incorrectly labeled if it was randomly labeled according to the distribution of labels in the subset:
\begin{equation}
    G(S) = 1 - \sum_{i=1}^{|C|} p_i^2
\end{equation}
The gradient of the Gini index is computationally cheaper to calculate because it avoids the logarithmic operations required by Shannon Entropy. 

\subsubsection{Gradient Boosting Optimization and Newton-Raphson Step}
Unlike Random Forests which rely on Bagging (Bootstrap Aggregating) to reduce variance by averaging parallel trees, XGBoost \parencite{chen2016xgboost} utilizes Boosting to reduce bias by sequentially training trees to correct the residual errors of preceding trees. The objective function $\mathcal{L}^{(t)}$ at iteration $t$ contains both a loss term $l$ and a regularization term $\Omega$:
\begin{equation}
    \mathcal{L}^{(t)} = \sum_{i=1}^n l(y_i, \hat{y}_i^{(t-1)} + f_t(x_i)) + \Omega(f_t)
\end{equation}
XGBoost optimizes this objective by taking a second-order Taylor expansion around the current estimate, functionally executing a Newton-Raphson step in function space:
\begin{equation}
    \mathcal{L}^{(t)} \simeq \sum_{i=1}^n \left[ l(y_i, \hat{y}_i^{(t-1)}) + g_i f_t(x_i) + \frac{1}{2} h_i f_t^2(x_i) \right] + \Omega(f_t)
\end{equation}
where $g_i$ and $h_i$ are the first and second order gradient statistics on the loss function, respectively. This rigorous mathematical foundation is exactly why XGBoost \parencite{chen2016xgboost} outperforms the Multi-Layer Perceptron \parencite{zhang2013review} on the sharp, threshold-based synthetic data generated for this thesis.

\subsection{Explainable AI (XAI): From LIME \parencite{lundberg2017unified} to Counterfactuals}
The deployment of machine learning in critical public sectors, such as education and career guidance, necessitates transparency. The European Union's General Data Protection Regulation (GDPR) explicitly mandates a "right to explanation" for algorithmic decisions. 

\subsubsection{Local Interpretable Model-agnostic Explanations (LIME)}
LIME operates by perturbing the input data and fitting a linear surrogate model around the local vicinity of the prediction. While LIME \parencite{lundberg2017unified} provides feature attributions, it fails to provide actionable pathways. A student informed that "Biology contributed 40\% to your rejection" gains no pedagogical recourse.

\subsubsection{Counterfactual Optimization Bounds}
Counterfactuals solve this by posing the question: "What is the minimum alteration required to flip the prediction?" Mathematically, this is an optimization problem seeking the perturbation vector $\delta$:
\begin{equation}
    \delta^* = \arg \min_{\delta} || \delta ||_p \quad \text{subject to} \quad f(x + \delta) = y_{target} \text{ and } x+\delta \in \mathcal{F}
\end{equation}
where $|| \cdot ||_p$ is a distance metric (commonly $L_1$ or $L_2$) and $\mathcal{F}$ represents the space of feasible actions (e.g., preventing negative changes, as a student cannot deliberately lower their past scores to achieve a better outcome).

% Expanded iteration 14 for extreme volume

